% -----------------------------------------------------------------------------
% common/lab-memo.tex
% ----------------------------------------------------------------------------- 
Instead of writing a detailed lab report each week, you will write a short memorandum, or \emph{memo}.  This is a common form of communication in business, science, and industry.

\begin{quote}
Memos are written by all ``knowledge workers.'' They are most often designed to be used within the writer's own company or organization.  They may note the existence of a problem, propose some course of action, describe a procedure, or report the results of a test or an investigation.  They are sometimes referred to as informal writing, which is not the same as sloppy, casual, or careless writing.  A memo must be carefully prepared, thoughtfully written, and thoroughly proofread for errors. \\
\hfill --- from \emph{A Survey of Industrial Mathematics} by Charles R.\ MacCluer.  Dover (2010).  Pages 282--283.
\end{quote}

A lab memo is a clear and concise report on your thoughts, work, observations, and conclusions.  Communication is the goal.  It should be an organized and coherent narrative of what you did and what you learned --- after you have had a chance to reflect and make sense of it.

You will write one lab memo each week.  When you finish gathering data and crunching the numbers, you should outline a short memo.  Your memo should be typed, and it is due 48 hours after your lab session ends. You should submit your memo as a PDF through \HomeworkSubmission.  You may include tables, graphs, and other notes in the Discussion.

Memos follow a standard format.  The one we will use for this lab is simple.

\newpage
\noindent\rule{\textwidth}{2px}
\begin{verbatim}
To:                 MEE@OIT Research Group
                    Oregon Institute of Technology
From:               Jesse Kinder
Subject:            [The topic of the memo]
Date:               [Today's Date]

Foreword:


Summary:

\newpage
Discussion:


\end{verbatim}
\noindent\rule{\textwidth}{2px}
\vfill

The body of your memo should contain three sections:
\begin{itemize}[itemsep=5pt]
\item\textbf{Foreword:}
In one or two sentences, describe the problem under investigation.

What was the goal of this work?  Provide a clear and concise overview.  Focus on the Big Picture.  Save the details for the summary.

\item\textbf{Summary:}
In one to three paragraphs, briefly describe your hypotheses, summarize your procedure, and report your observations.  Include numerical results --- with appropriate units and significant figures --- when relevant.

In general, the foreword and summary should fit on a single page.

\item\textbf{Discussion:}
Provide information that an interested reader might need to better understand your summary.

A good discussion includes at least one paragraph on at least one of the following.
\begin{itemize}[noitemsep]
	\item Summarize and interpret graphs, drawings, or calculations.
	\item Reflect on your work.  What did you find interesting?  Tedious?  Frustrating?
	\item Describe technical challenges.  What problems did you run into?
	\item Assess your work.  How confident are you in your numerical results?  How did you quantify uncertainty?
	\item Look ahead.  If we were to study this topic further, what would be the next steps?
	\item Describe special equipment or methods you used for the first time.
	\item Improve the work.  Is there a better way to do the experiments or the analyses?
\end{itemize}
\end{itemize}
